% ==================== 6.4 实验和结果分析 ====================
\section{实验和结果分析}
\label{sec:ch6_experiments}

为验证A-D3QN混合路径规划算法的有效性，本节设计了一系列实验，与基准算法进行对比分析。

\subsection{实验设置}
\label{subsec:ch6_setup}

\textbf{（1）实验环境}

\begin{itemize}
    \item 硬件：Intel Core i9-10900K CPU, NVIDIA RTX 3090 GPU, 64 GB内存
    \item 软件：Python 3.8, PyTorch 1.9, ROS Noetic
    \item 仿真：Gazebo + 月面场景模型
\end{itemize}

\textbf{（2）测试场景}

设计四种测试场景：
\begin{itemize}
    \item 场景A：平坦地形，稀疏障碍，尺寸100 m × 100 m；
    \item 场景B：起伏地形，中等障碍密度，尺寸200 m × 200 m；
    \item 场景C：复杂地形，密集障碍，尺寸300 m × 300 m；
    \item 场景D：动态障碍场景，障碍物随机移动。
\end{itemize}

\textbf{（3）对比算法}

\begin{itemize}
    \item A*算法：经典全局路径规划算法；
    \item RRT*算法：快速随机搜索树算法；
    \item DQN：标准深度Q网络；
    \item D3QN：Double Dueling DQN；
    \item A-D3QN：本文提出的算法。
\end{itemize}

\textbf{（4）评估指标}

\begin{itemize}
    \item 路径长度（m）：规划路径的总长度；
    \item 规划时间（ms）：算法运行时间；
    \item 成功率（\%）：成功到达目标的比例；
    \item 平滑度（\degree）：路径角度变化的均值；
    \item 能耗（J）：估计的能耗。
\end{itemize}

\subsection{算法验证与性能分析}
\label{subsec:ch6_results}

\textbf{（1）静态场景性能}

表\ref{tab:static_result}展示了静态场景下各算法的性能对比。

\begin{table}[htbp]
    \centering
    \caption{静态场景性能对比}
    \label{tab:static_result}
    \begin{tabular}{cccccc}
        \toprule
        场景 & 算法 & 路径长度（m） & 时间（ms） & 成功率 & 平滑度（\degree） \\
        \midrule
        \multirow{5}{*}{A} & A* & 142.5 & 12 & 100\% & 45.2 \\
         & RRT* & 156.3 & 85 & 95\% & 52.8 \\
         & DQN & 148.2 & 8 & 92\% & 38.5 \\
         & D3QN & 145.8 & 9 & 95\% & 35.2 \\
         & A-D3QN & \textbf{143.6} & 10 & \textbf{100\%} & \textbf{28.6} \\
        \midrule
        \multirow{5}{*}{C} & A* & 385.2 & 45 & 85\% & 48.5 \\
         & RRT* & 412.5 & 180 & 78\% & 55.2 \\
         & DQN & 398.6 & 15 & 88\% & 42.3 \\
         & D3QN & 392.8 & 18 & 91\% & 38.6 \\
         & A-D3QN & \textbf{378.5} & 20 & \textbf{97\%} & \textbf{32.1} \\
        \bottomrule
    \end{tabular}
\end{table}

结果表明，A-D3QN在复杂场景C中优势明显：路径长度比A*缩短1.7\%，成功率提升12个百分点，平滑度降低33.8\%。

\textbf{（2）动态场景性能}

表\ref{tab:dynamic_result}展示了动态场景D下的性能对比。

\begin{table}[htbp]
    \centering
    \caption{动态场景性能对比}
    \label{tab:dynamic_result}
    \begin{tabular}{cccc}
        \toprule
        算法 & 成功率 & 平均避障次数 & 平均时间（s） \\
        \midrule
        A*（重规划） & 72\% & 5.2 & 185.6 \\
        RRT* & 68\% & 6.8 & 210.3 \\
        DQN & 85\% & 4.5 & 168.2 \\
        D3QN & 89\% & 3.8 & 155.6 \\
        A-D3QN & \textbf{96\%} & \textbf{2.6} & \textbf{142.8} \\
        \bottomrule
    \end{tabular}
\end{table}

A-D3QN在动态场景中成功率最高（96\%），避障次数最少（2.6次），体现了良好的适应性和决策能力。

\textbf{（3）训练收敛性}

图\ref{fig:training}展示了各算法的训练收敛曲线。A-D3QN在约800个episode后收敛，比DQN快约20\%。

\begin{figure}[htbp]
    \centering
    \begin{tikzpicture}
        \begin{axis}[
            width=10cm, height=5cm,
            xlabel={训练轮次（Episode）},
            ylabel={平均奖励},
            xmin=0, xmax=1500,
            ymin=-100, ymax=200,
            legend pos=south east,
            grid=major
        ]
            \addplot[blue, thick] coordinates {
                (0, -80) (200, -20) (400, 50) (600, 100) (800, 140) (1000, 160) (1200, 170) (1500, 175)
            };
            \addplot[red, thick] coordinates {
                (0, -85) (200, -30) (400, 30) (600, 80) (800, 120) (1000, 140) (1200, 150) (1500, 155)
            };
            \addplot[green, thick] coordinates {
                (0, -90) (200, -40) (400, 10) (600, 60) (800, 100) (1000, 120) (1200, 130) (1500, 135)
            };
            \legend{A-D3QN, D3QN, DQN}
        \end{axis}
    \end{tikzpicture}
    \caption{训练收敛曲线}
    \label{fig:training}
\end{figure}

\textbf{（4）能耗分析}

表\ref{tab:energy}展示了各算法生成的路径能耗对比。

\begin{table}[htbp]
    \centering
    \caption{路径能耗对比}
    \label{tab:energy}
    \begin{tabular}{ccc}
        \toprule
        算法 & 能耗（kJ） & 相对节省 \\
        \midrule
        A* & 125.6 & 基准 \\
        DQN & 118.2 & 5.9\% \\
        D3QN & 112.8 & 10.2\% \\
        A-D3QN & \textbf{105.3} & \textbf{16.2\%} \\
        \bottomrule
    \end{tabular}
\end{table}

A-D3QN生成的路径能耗最低，比A*节省16.2\%，原因是算法考虑了坡度、粗糙度等因素，选择了更节能的路径。

\textbf{（5）消融实验}

表\ref{tab:ablation}展示了A-D3QN各组件的消融实验结果。

\begin{table}[htbp]
    \centering
    \caption{消融实验结果}
    \label{tab:ablation}
    \begin{tabular}{lcc}
        \toprule
        模型配置 & 成功率 & 相对提升 \\
        \midrule
        基准DQN & 88\% & 基准 \\
        + Double Q & 91\% & +3\% \\
        + Dueling结构 & 93\% & +5\% \\
        + 注意力机制 & \textbf{97\%} & +9\% \\
        \bottomrule
    \end{tabular}
\end{table}

注意力机制的引入带来最大提升（+4\%），证明了注意力机制在路径规划任务中的有效性。

综合实验结果表明，A-D3QN混合路径规划算法在静态和动态场景下均表现出色，生成的路径更短、更平滑、更节能，满足月面巡视器自主行走的实际需求。
